% ============================================================
% CHAPTER 20: MULTICUT IN GENERAL GRAPHS (Vazirani Ch. 20)
% ============================================================
\setcounter{chapter}{19}
\chapter{Multicut in General Graphs}
\label{ch:multicut_general}

\section{Intuition: Region Growing and LP Rounding}

\begin{intuitionbox}
\textbf{Region Growing:} The Multicut problem asks for a minimum cost set of edges whose removal disconnects given pairs of vertices. While NP-hard in general graphs, the problem admits an $O(\log k)$ approximation via LP rounding. The key idea is \emph{region growing}: starting from each source $s_i$, grow a ball using edge costs weighted by $1/x_e^*$ (the reciprocal of the LP solution). When the ball reaches $t_i$, we cut the boundary edges. The region growing ensures that each ball's boundary is cheap relative to the LP value.
\end{intuitionbox}

\section{Problem Definition}
Given an undirected graph $G = (V, E)$ with edge costs $c_e \ge 0$, and $k$ demand pairs $(s_1, t_1), \dots, (s_k, t_k)$, find a minimum cost subset $C \subseteq E$ such that $s_i$ and $t_i$ are in different connected components of $(V, E \setminus C)$ for all $i$.

The problem is NP-hard in general graphs. It admits an $O(\log k)$ approximation.

\section{LP Relaxation}
\[
\min \sum_{e \in E} c_e x_e \quad \text{s.t.} \quad \sum_{e \in P_i} x_e \ge 1 \; (\forall i), \quad x_e \ge 0
\]
where $P_i$ is the set of edges on any path from $s_i$ to $t_i$.

\section{Algorithm and Python Implementation}
The algorithm solves the LP, then uses region growing: for each pair $(s_i, t_i)$, grow a ball from $s_i$ using Dijkstra with edge weights $c_e / x_e^*$. When $t_i$ is reached, cut all boundary edges of the ball.

\begin{lstlisting}[style=python, caption=Multicut in General Graphs — Region Growing]
import heapq
from lp_algorithms import Simplex   # Simplex from Chapter 12

def solve_multicut_lp(n, edges, costs, pairs):
    """Fractional x* for Multicut via the dual LP
       (max sum y_j  st. sum_{j: e in path_j} y_j <= c_e, y_j >= 0)."""
    adj = {v: [] for v in range(n)}
    for u, v in edges:
        adj[u].append(v); adj[v].append(u)
    eidx = {tuple(sorted(e)): i for i, e in enumerate(edges)}

    # Enumerate simple s_i-t_i paths for each demand pair
    paths = []
    def dfs(u, t, visited, path):
        if u == t:
            paths.append(list(path)); return
        for v in adj[u]:
            if v not in visited:
                visited.add(v); path.append(tuple(sorted((u, v))))
                dfs(v, t, visited, path)
                path.pop(); visited.remove(v)

    for s, t in pairs:
        dfs(s, t, {s}, [])

    if not paths:
        return [0.0] * len(edges), 0.0

    # Dual: max 1^T y   s.t.  A^T y <= costs,  y >= 0
    A = [[0.0] * len(paths) for _ in edges]
    for j, p in enumerate(paths):
        for e in p:
            A[eidx[e]][j] = 1.0
    solver = Simplex(A, costs[:], [1.0] * len(paths))
    y, dual_obj = solver.solve()
    x = [max(0.0, min(1.0, v)) for v in solver.obj_row[len(paths):len(paths)+len(edges)]]
    return x, dual_obj

def multicut_region_growing(n, edges, costs, pairs):
    """O(log k) approximation for Multicut via LP + region growing."""
    x, _ = solve_multicut_lp(n, edges, costs, pairs)
    adj = {v: [] for v in range(n)}
    for u, v in edges:
        adj[u].append(v); adj[v].append(u)
    eidx = {tuple(sorted(e)): i for i, e in enumerate(edges)}
    cut = set()
    INF = float('inf')

    for s, t in pairs:
        dist = [INF] * n; dist[s] = 0.0
        visited = [False] * n
        pq = [(0.0, s)]
        while pq:
            d, u = heapq.heappop(pq)
            if visited[u]: continue
            visited[u] = True
            if u == t: break
            for v in adj[u]:
                if not visited[v]:
                    e = tuple(sorted((u, v)))
                    w = costs[eidx[e]] / max(x[eidx[e]], 1e-9)
                    if d + w < dist[v]:
                        dist[v] = d + w; heapq.heappush(pq, (dist[v], v))

        ball = {v for v in range(n) if dist[v] < dist[t] - 1e-9}
        for u, v in edges:
            if (u in ball) != (v in ball):
                cut.add(tuple(sorted((u, v))))

    cut_edges = list(cut)
    total = sum(costs[eidx[e]] for e in cut_edges)
    return cut_edges, total
\end{lstlisting}

\section{Analysis: $O(\log k)$ Approximation}
Each ball has boundary cost at most $O(\log k)$ times the LP cost contribution for that pair. Summing over all pairs gives the $O(\log k)$ approximation ratio. The key lemma is that for any ball $B$ with boundary $\delta(B)$, the cost $\sum_{e \in \delta(B)} c_e \le O(\log k) \cdot \sum_{e \in \delta(B)} x_e^*$.

\section{Applications}
\begin{itemize}
    \item \textbf{Network Security:} Isolating sensitive communication pairs by cutting edges.
    \item \textbf{VLSI Design:} Separating signal nets on a chip.
    \item \textbf{Image Segmentation:} Partitioning pixels into regions by cutting edges in a graph.
\end{itemize}

\section{Exercises}
\begin{exercise}
Show that the LP relaxation for Multicut has an integrality gap of $\Omega(\log k)$.
\end{exercise}
\begin{exercise}
Implement a greedy alternative: repeatedly find the cheapest edge whose removal disconnects at least one pair that is still connected.
\end{exercise}

